\documentclass[11pt]{article}

\usepackage[margin=1.1in]{geometry}
\usepackage{amsmath,amssymb,amsthm}
\usepackage{booktabs}
\usepackage{listings}
\usepackage{xcolor}
\usepackage{microtype}
\emergencystretch=1.5em
\usepackage[hidelinks]{hyperref}

\newtheorem{theorem}{Theorem}[section]
\newtheorem{proposition}[theorem]{Proposition}
\newtheorem{lemma}[theorem]{Lemma}
\newtheorem{corollary}[theorem]{Corollary}
\theoremstyle{remark}
\newtheorem{remark}[theorem]{Remark}
\theoremstyle{definition}
\newtheorem{definition}[theorem]{Definition}

\lstdefinelanguage{lean}{
  keywords={theorem,def,noncomputable,by,fun,match,with,open,import,namespace,end},
  sensitive=true,
  comment=[l]{--},
  morecomment=[s]{/-}{-/},
}
\lstset{
  language=lean,
  basicstyle=\ttfamily\small,
  keywordstyle=\bfseries,
  commentstyle=\itshape\color{gray},
  columns=fullflexible,
  keepspaces=true,
  breaklines=true,
  literate={∑}{{$\textstyle\sum$}}1
           {∏}{{$\textstyle\prod$}}1
           {∈}{{$\in$}}1
           {⊤}{{$\top$}}1
           {ℕ}{{$\mathbb{N}$}}1
           {ℝ}{{$\mathbb{R}$}}1
           {⁻¹}{{$^{-1}$}}2
           {→}{{$\to$}}1
           {↔}{{$\leftrightarrow$}}1
           {∀}{{$\forall$}}1
           {≠}{{$\neq$}}1,
}

\title{A Machine-Verified Bijective Proof of the Rooted\\
Child-Factorial Catalan Identity over Spanning Trees\\
of the Complete Graph}
\author{Lluis Eriksson}
\date{July 2, 2026}

\begin{document}
\maketitle

\begin{abstract}
Let $K_{n+1}$ be the complete graph on the vertex set $\{0,1,\dots,n\}$, and
for a spanning tree $T$ of $K_{n+1}$, rooted at $0$, let $c_T(v)$ denote the
number of children of the vertex $v$.  We prove the exact identity
\[
  \sum_{T \in \mathrm{ST}(K_{n+1})} \; \prod_{v=0}^{n} c_T(v)! \;=\; n!\,C_n ,
\]
where $C_n$ is the $n$-th Catalan number; equivalently, the normalized sum
$(n+1)\,\bigl((n+1)!\bigr)^{-1}\sum_T \prod_v c_T(v)!$ equals $C_n$ exactly.
The proof is bijective: pairs consisting of a spanning tree together with a
linear ordering of every child set are placed in explicit bijection with
vertex-labeled plane trees on $n+1$ nodes whose root carries the label $0$.
The identity arises as the exact ``second-Ursell'' normalization constant in
the author's audit-first programme on four-dimensional $SU(N)$ Yang--Mills
existence and mass gap, where it had been isolated as a named open proposition
in a public challenge repository; the present paper is self-contained
combinatorics and makes no claim about that programme.  The entire proof has
been formalized in Lean~4 against a pinned Mathlib snapshot: the headline
declarations compile with no \texttt{sorry}, and the kernel's axiom oracle
reports exactly \texttt{[propext, Classical.choice, Quot.sound]}.  All
artifacts, including a pinned continuous-integration replay of the full
verification, are public.
\end{abstract}

\section{Introduction}\label{sec:intro}

Cluster and tree-graph expansions bound logarithms of partition functions by
sums over trees; this line of reasoning goes back to Penrose's tree-graph
identity~\cite{Penrose1967} and was given its modern abstract-polymer form by
Koteck\'y and Preiss~\cite{KoteckyPreiss1986}.  In quantitative applications
one repeatedly meets weighted sums over the spanning trees of a complete
graph, where the weight of a tree is a product of local factors attached to
its vertices.  The simplest nontrivial member of this family assigns to a
rooted spanning tree the product of the factorials of its child counts.

In the author's programme on the construction of four-dimensional $SU(N)$
Yang--Mills theory~\cite{Eriksson2026Master,Eriksson2026KP}, the terminal
Koteck\'y--Preiss estimate requires the \emph{exact} evaluation of this
factorially weighted spanning-tree sum: the ``second-Ursell'' normalization
constant of the expansion is precisely the ratio of the weighted sum to
$(n+1)!/(n+1) = n!$, and convergence at the stated radius needs this ratio to
be the Catalan number $C_n$ on the nose, not merely of Catalan order.  The
statement was therefore isolated, under the name
\texttt{RootedChildFactorialCatalanIdentity}, as an explicitly open
proposition in a public challenge repository, with the request that it be
closed by an unconditional, machine-checked proof against the repository's own
definitions.

This paper does three things.
\begin{enumerate}
  \item It states and proves the identity
    (Theorem~\ref{thm:main}) by an explicit bijection between pairs
    \emph{(spanning tree, orderings of all child sets)} and vertex-labeled
    plane trees with root label $0$ (Section~\ref{sec:bijection}), combined
    with the enumeration of the latter objects
    (Section~\ref{sec:counting}).
  \item It reports a complete formalization of the proof in
    Lean~4~\cite{Lean4} over the Mathlib library~\cite{mathlib}, carried out
    against the pinned toolchain of the challenge repository and checked by
    the Lean kernel with a clean axiom footprint
    (Section~\ref{sec:lean}).
  \item It documents the reproducibility artifacts: the repository, the
    pull request closing the challenge, the exact commit, and a pinned
    continuous-integration replay of the entire verification
    (Section~\ref{sec:repro}).
\end{enumerate}
The mathematics below is entirely self-contained and elementary; no result
from, and no claim about, the Yang--Mills programme is used or made.

\section{Statement and numerical audit}\label{sec:statement}

Fix $n \ge 1$ and let $V = \{0,1,\dots,n\}$, so that $|V| = n+1$.  We write
$\mathrm{ST}(K_{n+1})$ for the set of spanning trees of the complete graph on
$V$, encoded as sets of unordered edges.  Every spanning tree $T$ is rooted at
the distinguished vertex $0$: each $w \neq 0$ has a unique neighbor lying on
the path from $w$ to $0$, its \emph{parent} $p_T(w)$, and
\[
  \mathrm{ch}_T(v) \;=\; \{\, w \neq 0 : p_T(w) = v \,\}, \qquad
  c_T(v) \;=\; |\mathrm{ch}_T(v)|
\]
are the \emph{children} and the \emph{child count} of $v$.  (In the
formalization the parent is defined for an arbitrary edge set as the minimal
neighbor one breadth-first-search level closer to the root; for a tree this
minimum is over a singleton and coincides with the graph-theoretic parent.
All statements below are proved against that general definition.)

Let $C_n = \frac{1}{n+1}\binom{2n}{n}$ denote the Catalan numbers, so
$C_1, C_2, C_3, \dots = 1, 2, 5, 14, 42, 132, \dots$

\begin{theorem}[rooted child-factorial Catalan identity]\label{thm:main}
For every $n \ge 1$,
\[
  \sum_{T \in \mathrm{ST}(K_{n+1})} \; \prod_{v \in V} c_T(v)!
  \;=\; n! \, C_n .
\]
\end{theorem}

\begin{corollary}[normalized form]\label{cor:normalized}
For every $n \ge 1$,
\[
  (n+1)\,\bigl((n+1)!\bigr)^{-1}
  \sum_{T \in \mathrm{ST}(K_{n+1})} \; \prod_{v \in V} c_T(v)!
  \;=\; C_n
\]
as an identity of real numbers.  This is the form
\emph{\texttt{RootedChildFactorialCatalanIdentity}} of the challenge
repository.
\end{corollary}

\begin{remark}
The unweighted analogue of the left-hand side is Cayley's formula
$|\mathrm{ST}(K_{n+1})| = (n+1)^{n-1}$~\cite{Cayley1889}.  Since every weight
$\prod_v c_T(v)!$ is a positive integer, Theorem~\ref{thm:main} in particular
interpolates between the two classical sequences: for $n = 4$, for instance,
the $125$ spanning trees of $K_5$ carry total weight $336 = 4! \cdot 14$.
\end{remark}

Before turning to the proof we record a brute-force audit.  Enumerating all
spanning trees of $K_{n+1}$ through Pr\"ufer sequences and computing the
weighted sum directly gives the following table, in exact agreement with
Theorem~\ref{thm:main}.

\begin{center}
\begin{tabular}{rrrr}
\toprule
$n$ & $|\mathrm{ST}(K_{n+1})| = (n+1)^{n-1}$ &
$\sum_T \prod_v c_T(v)!$ & $n!\,C_n$ \\
\midrule
1 & 1 & 1 & 1 \\
2 & 3 & 4 & 4 \\
3 & 16 & 30 & 30 \\
4 & 125 & 336 & 336 \\
5 & 1296 & 5040 & 5040 \\
6 & 16807 & 95040 & 95040 \\
\bottomrule
\end{tabular}
\end{center}

\section{The bijection}\label{sec:bijection}

The left-hand side of Theorem~\ref{thm:main} counts, with multiplicity, the
pairs
\[
  (T, L), \qquad T \in \mathrm{ST}(K_{n+1}), \quad
  L = (\ell_v)_{v \in V},
\]
where $\ell_v$ is a duplicate-free list enumerating the child set
$\mathrm{ch}_T(v)$; for fixed $T$ there are exactly $\prod_v c_T(v)!$ such
\emph{child listings}.  We call the pairs $(T,L)$ \emph{ordered rooted
spanning trees}.

The right-hand side counts \emph{labeled plane trees}.  A plane tree (rose
tree) with labels in $V$ is either a node carrying a label $a \in V$ together
with a finite \emph{ordered} list of labeled plane subtrees.  Write
$\mathrm{lab}(u)$ for the list of labels of $u$ in root-first order, and let
\[
  \mathcal{L}_n \;=\; \{\, u :
    \mathrm{lab}(u) \text{ is a permutation of } (0,1,\dots,n)
    \text{ and the root label of } u \text{ is } 0 \,\}.
\]
Section~\ref{sec:counting} shows $|\mathcal{L}_n| = n!\,C_n$.
Theorem~\ref{thm:main} therefore follows from:

\begin{theorem}[bijection]\label{thm:bij}
The map
\[
  \Phi : (T, L) \;\longmapsto\; \Phi_{T,L}(0),
\]
where $\Phi_{T,L}(v)$ is the labeled plane tree with root label $v$ whose
ordered subtrees are $\Phi_{T,L}(w)$ for $w$ running through $\ell_v$ in list
order, is a bijection from the set of ordered rooted spanning trees onto
$\mathcal{L}_n$.
\end{theorem}

The recursion defining $\Phi_{T,L}$ is well founded because a child sits one
breadth-first-search level below its parent, so the quantity
$(n+1) - \mathrm{level}_T(v)$ strictly decreases along the recursion.  The
proof of Theorem~\ref{thm:bij} splits into the three usual claims, whose
formalized statements are quoted in Section~\ref{sec:lean}.

\subsection{Well-definedness}

\begin{proposition}\label{prop:welldef}
For every ordered rooted spanning tree $(T,L)$ one has
$\Phi_{T,L}(0) \in \mathcal{L}_n$.
\end{proposition}

\begin{proof}[Proof sketch]
The label multiset of the subtree $\Phi_{T,L}(v)$ is characterized by parent
iterates:
\[
  w \in \mathrm{lab}\bigl(\Phi_{T,L}(v)\bigr)
  \iff
  \exists k \ge 0 : \; p_T^{\,k}(w) = v
  \text{ and } \mathrm{level}_T(w) = \mathrm{level}_T(v) + k .
\]
(The proof is by the same well-founded recursion as the definition of
$\Phi$.)  Two consequences follow.  First, subtrees rooted at distinct
children of the same node have disjoint label sets: a common label $w$ would
have the two children as parent iterates at the \emph{same} offset
$k = \mathrm{level}(w) - \mathrm{level}(\text{child})$, hence the children
would coincide.  This makes $\mathrm{lab}(\Phi_{T,L}(v))$ duplicate-free, by
induction.  Second, at $v = 0$ every vertex $w$ occurs, with
$k = \mathrm{level}_T(w)$, because iterating the parent map from $w$ reaches
the root in exactly $\mathrm{level}_T(w)$ steps.  A duplicate-free list of
$n+1$ elements of $V$ is a permutation of $(0,\dots,n)$, and the root label of
$\Phi_{T,L}(0)$ is $0$ by construction.
\end{proof}

\subsection{Injectivity}

For a labeled plane tree $u$ with duplicate-free labels and $v \in
\mathrm{lab}(u)$, let $\mathrm{childrenOf}(u,v)$ denote the list of root
labels of the ordered subtrees of the unique node of $u$ labeled $v$.

\begin{proposition}\label{prop:inj}
$\Phi$ is injective.  More precisely, if $u = \Phi_{T,L}(0)$ then
$\ell_v = \mathrm{childrenOf}(u,v)$ for every $v$, and $T$ is recovered as the
set of parent edges $\{\{v, w\} : w \in \ell_v\}$.
\end{proposition}

\begin{proof}[Proof sketch]
The extraction $\ell_v = \mathrm{childrenOf}(u,v)$ is proved by the same
well-founded recursion, using the pairwise disjointness of sibling label sets
to show that the search for the node labeled $v$ descends into the correct
subtree.  Thus $L$ is determined by $u$.  Next, the listings determine the
parent map ($w \in \ell_v$ forces $p_T(w) = v$), and a spanning tree is
determined by its parent map: its edge set \emph{equals} its set of parent
edges.  The latter recovery statement is a lemma of the upstream artifact (in
the terminology of the repository, $\mathrm{penroseTree}(T) = T$ for every
spanning tree $T$), and injectivity follows.
\end{proof}

\subsection{Surjectivity}

\begin{proposition}\label{prop:surj}
$\Phi$ is surjective onto $\mathcal{L}_n$.
\end{proposition}

\begin{proof}[Proof sketch]
Given $u \in \mathcal{L}_n$, define
\[
  T_u = \bigl\{ \{v, w\} : w \in \mathrm{childrenOf}(u,v) \bigr\},
  \qquad
  L_u = \bigl(\mathrm{childrenOf}(u,v)\bigr)_{v \in V}.
\]
Three facts are needed.

\emph{(a) $T_u$ is a spanning tree.}  Connectivity: by structural induction
on $u$, every label reaches the root label along parent--child edges.  Edge
count: the edges of $T_u$ are in bijection with the disjoint union
$\bigsqcup_v \mathrm{childrenOf}(u,v)$ (distinct parents contribute distinct
edges, by antisymmetry of the structural parent--child relation on a
duplicate-free tree), and that disjoint union is in bijection with the set of
non-root labels (every non-root label has a structural parent, and the parent
is unique).  Hence $|T_u| = n$, and a connected graph on $n+1$ vertices with
$n$ edges is a tree.

\emph{(b) The rooted children of $T_u$ are the structural children:}
$\mathrm{ch}_{T_u}(v) = \mathrm{childrenOf}(u,v)$ as sets.  The key claim is
that for every $w \neq 0$,
\[
  w \in \mathrm{childrenOf}\bigl(u,\, p_{T_u}(w)\bigr),
\]
proved by strong induction on the breadth-first-search level of $w$ in $T_u$.
Indeed $p_{T_u}(w)$ is adjacent to $w$, and adjacency in $T_u$ means exactly
``structural parent or structural child''.  If $b := p_{T_u}(w)$ were a
structural \emph{child} of $w$, then $b \neq 0$ (the root label $0$ is never a
child), the inductive hypothesis applies to $b$ (its level is smaller by the
parent equation), and by uniqueness of structural parents $p_{T_u}(b) = w$.
The two level equations
$\mathrm{level}(b) + 1 = \mathrm{level}(w)$ and
$\mathrm{level}(w) + 1 = \mathrm{level}(b)$
then contradict each other.  Notably, this argument needs neither a depth
function on $u$ nor an explicit parent function: the two instances of the
parent specification collide by themselves.

\emph{(c) Reconstruction:} $\Phi_{T_u, L_u}(0) = u$.  This is an instance of
the general lemma that $\Phi_{T,L}(\mathrm{root}(u)) = u$ whenever the
listings $L$ agree with $\mathrm{childrenOf}(u,\cdot)$ on the labels of $u$
--- proved by structural induction on $u$, the agreement hypothesis
propagating to each subtree because $\mathrm{childrenOf}$ restricted to the
labels of a subtree coincides with the subtree's own $\mathrm{childrenOf}$.
For $(T_u, L_u)$ the agreement holds by definition ($L_u$ \emph{is}
$\mathrm{childrenOf}(u,\cdot)$), and $\mathrm{root}(u) = 0$.
\end{proof}

\section{Counting labeled plane trees}\label{sec:counting}

\begin{proposition}\label{prop:count}
$|\mathcal{L}_n| = n! \, C_n$ for every $n \ge 1$.
\end{proposition}

\begin{proof}[Proof sketch]
Unlabeled plane trees with $m+1$ nodes are counted by $C_m$: deleting the
root's first subtree decomposes a plane tree with $m+1$ nodes into an ordered
pair of strictly smaller plane trees, which yields the recurrence
$P_{m+1} = \sum_{i+j = m} P_i\,P_j$ with $P_0 = 1$ (nodes counted so that
$P_m$ is the number of plane trees with $m+1$ nodes), i.e.\ exactly the
defining recurrence of the Catalan numbers, as implemented in Mathlib's
\texttt{Combinatorics.Enumerative.Catalan}.  A labeled plane tree in
$\mathcal{L}_n$ is equivalently a pair (shape, labeling): writing the labels
in root-first order identifies the labelings of a fixed shape with the
permutations of $(0,\dots,n)$ whose first entry is $0$, of which there are
$n!$, and the pair (shape, label word) determines the labeled tree uniquely
(rigidity of the decoration).  Hence
$|\mathcal{L}_n| = n! \cdot C_n$.
\end{proof}

\begin{proof}[Proof of Theorem~\ref{thm:main} and
Corollary~\ref{cor:normalized}]
By Theorem~\ref{thm:bij} and Proposition~\ref{prop:count},
\[
  \sum_{T} \prod_v c_T(v)!
  \;=\; \#\{(T,L)\}
  \;=\; |\mathcal{L}_n|
  \;=\; n!\,C_n .
\]
For the corollary, multiply by $(n+1)\bigl((n+1)!\bigr)^{-1} = (n!)^{-1}$ and
cancel: $(n+1)\,\bigl((n+1)!\bigr)^{-1}\, n!\, C_n = C_n$, using
$(n+1)! = (n+1)\,n!$ and $n! \neq 0$.
\end{proof}

\section{The Lean formalization}\label{sec:lean}

Every statement above has been formalized in Lean~4~\cite{Lean4} over
Mathlib~\cite{mathlib}, against the challenge repository's pinned toolchain
(Lean \texttt{leanprover/lean4:v4.29.0-rc6}, Mathlib commit
\texttt{07642720480157414db592fa85b626dafb71355b}) and against the
repository's \emph{own} definitions of spanning trees, breadth-first-search
parents, and rooted child counts.  The headline declarations are:

\begin{lstlisting}
theorem sum_prod_rootedChildCount_factorial_eq (n : ℕ) :
    (∑ T ∈ spanningTrees (⊤ : SimpleGraph (Fin (n + 1))),
        ∏ v, (rootedChildCount T v)!) = n ! * catalan n

theorem rootedChildFactorialCatalanIdentity_holds (n : ℕ) :
    RootedChildFactorialCatalanIdentity n
\end{lstlisting}

\noindent
where the second statement unfolds, per the repository's definitions, to the
real-number identity of Corollary~\ref{cor:normalized}:

\begin{lstlisting}
noncomputable def rootedChildFactorialTreeSum (n : ℕ) : ℝ :=
  ∑ T ∈ spanningTrees (⊤ : SimpleGraph (Fin (n + 1))),
    ∏ v : Fin (n + 1), ((rootedChildCount T v).factorial : ℝ)

noncomputable def rootedChildFactorialTreeSumNormalized (n : ℕ) : ℝ :=
  ((n : ℝ) + 1) * (((n + 1).factorial : ℝ))⁻¹ *
    rootedChildFactorialTreeSum n
\end{lstlisting}

The proof development comprises five modules
(about $2{,}200$ lines in total):

\begin{center}
\begin{tabular}{lrl}
\toprule
Module & Lines & Role \\
\midrule
\texttt{RootedCatalanWord} & 234 &
  word-side algebra; normalized-form arithmetic \\
\texttt{PlaneTreeCatalan} & 214 &
  plane trees with $m+1$ nodes are counted by $C_m$ \\
\texttt{LabeledPlaneTreeCatalan} & 382 &
  decoration rigidity; $|\mathcal{L}_n| = n!\,C_n$ \\
\texttt{RootedCatalanBridge} & 1288 &
  the bijection: Props.~\ref{prop:welldef}--\ref{prop:surj},
  Thm.~\ref{thm:bij}, Thm.~\ref{thm:main} \\
\texttt{RootedCatalanIdentityProof} & 61 &
  Corollary~\ref{cor:normalized} against the repository's definitions \\
\bottomrule
\end{tabular}
\end{center}

A few engineering points may be of independent interest to formalizers.

\begin{itemize}
  \item \emph{Well-founded tree building.}  The map $\Phi$ is defined by
    well-founded recursion on $(n+1) - \mathrm{level}_T(v)$; the label
    characterization of Proposition~\ref{prop:welldef} and the extraction of
    Proposition~\ref{prop:inj} are proved by the same measure.
  \item \emph{Nested induction on rose trees.}  Labeled plane trees are a
    nested inductive type (a node carries a \emph{list} of subtrees).  All
    structural lemmas --- uniqueness and existence of structural parents,
    antisymmetry of the parent--child relation, duplicate-freeness of child
    listings, reachability of the root --- are proved by recursion on
    \texttt{sizeOf}, descending into list members via
    \texttt{List.sizeOf\_lt\_of\_mem}.
  \item \emph{The level-clash argument.}  Step (b) of
    Proposition~\ref{prop:surj} is formalized exactly as sketched: a strong
    induction on the breadth-first-search level in which the ``bad'' branch is
    refuted by instantiating the parent specification twice and letting the
    two resulting level equations contradict each other by linear arithmetic.
    No depth function on the plane tree is ever introduced.
  \item \emph{Counting by double disjoint union.}  The edge count
    $|T_u| = n$ is two applications of
    \texttt{Finset.card\_biUnion}: edges $\leftrightarrow$ child slots
    $\leftrightarrow$ non-root labels, with the required disjointness supplied
    by antisymmetry and by uniqueness of parents, respectively.
  \item \emph{Axiom footprint.}  The development contains no \texttt{sorry}
    and no additional axioms; for every exported declaration, including both
    headline theorems, \texttt{\#print axioms} reports exactly
\begin{verbatim}
[propext, Classical.choice, Quot.sound]
\end{verbatim}
    i.e.\ the standard classical axioms of Mathlib.
\end{itemize}

\section{Reproducibility and artifacts}\label{sec:repro}

All artifacts are public in the challenge repository
\begin{center}
\url{https://github.com/lluiseriksson/rooted-tree-catalan-closure}
\end{center}
The closure is submitted as pull request \#5 (branch
\texttt{codex/catalan-rigidity-lean-modules}, head commit
\texttt{1eff8fc651ed3338800b8ff67471e6091ce8ef33}), which adds the five proof
modules under \texttt{lean-patch/YangMills/KP/}, updates the theorem manifest
and claim-boundary documents, and archives the verification logs.

The verification is replayed from scratch by a pinned GitHub Actions
workflow; the run
\begin{center}
\url{https://github.com/lluiseriksson/rooted-tree-catalan-closure/actions/runs/28610948452}
\end{center}
(i) checks out the pinned upstream artifact and Mathlib snapshot,
(ii) applies the documented import-narrowing patch --- which leaves every
upstream proof body byte-identical and only trims module import blocks so the
relevant dependency chain builds in isolation ---
(iii) compiles the upstream adapter chain and the five new modules,
(iv) runs the axiom-oracle checks and prints \texttt{\#print axioms} for the
closing declarations, and
(v) validates the resulting logs against the archived evidence with the
repository's replay validators.
The archived logs are
\texttt{catalan-bridge-modules-replay.log},
\texttt{catalan-bridge-full-compile.log}, and
\texttt{catalan-identity-proof-oracle.log} under
\texttt{lean-patch/verification/}.

Independently of Lean, the brute-force audit of
Section~\ref{sec:statement} (a $30$-line Python script enumerating spanning
trees via Pr\"ufer sequences) confirms the identity for $1 \le n \le 6$.

\section{Related work and outlook}\label{sec:related}

The identity sits at the intersection of two classical themes.  On the
enumerative side, plane trees and their labelings belong to the standard
Catalan circle of ideas~\cite{Stanley2015}; the specific weighted sum over
spanning trees of $K_{n+1}$, with weight $\prod_v c_T(v)!$, appears to be less
standard, and we are not aware of a prior published bijective treatment,
although the result itself is elementary once the correct bijection is
identified.  On the analytic side, factorially weighted tree sums are the
combinatorial core of tree-graph and cluster
expansions~\cite{Penrose1967,KoteckyPreiss1986}; within the author's
programme, Theorem~\ref{thm:main} supplies the exact constant needed by the
terminal Koteck\'y--Preiss bound~\cite{Eriksson2026KP}, and the machine-checked
closure documented here continues the programme's mechanical-audit
methodology~\cite{Eriksson2026Audit}.

Two natural directions remain.  First, the same bijection should evaluate
refinements of the sum, e.g.\ with weights $\prod_v x_{c_T(v)}$ for formal
variables $x_k$, connecting to the exponential-generating-function calculus of
rooted labeled trees.  Second, on the formal side, the bridge module develops
a small reusable theory of labeled rose trees (structural parents, listings,
extraction, reconstruction) that may be worth contributing to Mathlib in
generic form.

\subsection*{Acknowledgments}
The Lean formalization was developed interactively with Anthropic's Claude;
every theorem stated here was checked by the Lean~4 kernel against the pinned
Mathlib snapshot, with the axiom oracle output archived in the repository.

\begin{thebibliography}{99}

\bibitem{Cayley1889}
A.~Cayley,
\emph{A theorem on trees},
Quart.\ J.\ Pure Appl.\ Math.\ \textbf{23} (1889), 376--378.

\bibitem{Penrose1967}
O.~Penrose,
\emph{Convergence of fugacity expansions for classical systems},
in: T.~A.~Bak (ed.), \emph{Statistical Mechanics: Foundations and
Applications}, Benjamin, New York, 1967, 101--109.

\bibitem{KoteckyPreiss1986}
R.~Koteck\'y and D.~Preiss,
\emph{Cluster expansion for abstract polymer models},
Comm.\ Math.\ Phys.\ \textbf{103} (1986), 491--498.

\bibitem{Stanley2015}
R.~P.~Stanley,
\emph{Catalan Numbers},
Cambridge University Press, 2015.

\bibitem{Lean4}
L.~de~Moura and S.~Ullrich,
\emph{The Lean 4 theorem prover and programming language},
in: Automated Deduction -- CADE 28, Lecture Notes in Computer Science
12699, Springer, 2021, 625--635.

\bibitem{mathlib}
The mathlib Community,
\emph{The Lean mathematical library},
in: Proceedings of the 9th ACM SIGPLAN International Conference on Certified
Programs and Proofs (CPP 2020), ACM, 2020, 367--381.

\bibitem{Eriksson2026Master}
L.~Eriksson,
\emph{The Master Map: An Audit-First, Attack-Resistant Navigation Guide to the
Unconditional Solution of the 4D $SU(N)$ Yang--Mills Existence and Mass Gap
Problem},
ai.viXra:2602.0096, 2026.

\bibitem{Eriksson2026KP}
L.~Eriksson,
\emph{Closing the Last Gap in the 4D $SU(N)$ Yang--Mills Construction: A
Verified Terminal KP Bound and an Explicit Clay Checklist},
ai.viXra:2602.0091, 2026.

\bibitem{Eriksson2026Audit}
L.~Eriksson,
\emph{Mechanical Audit Experiments and Reproducibility Appendix for a
Companion-Paper Programme on 4D $SU(N)$ Yang--Mills Existence and Mass Gap},
ai.viXra:2602.0117, 2026.

\bibitem{Repo}
L.~Eriksson,
\emph{rooted-tree-catalan-closure} (challenge repository),
\url{https://github.com/lluiseriksson/rooted-tree-catalan-closure},
pull request \#5, commit \texttt{1eff8fc6}, CI run 28610948452, 2026.

\end{thebibliography}

\end{document}
